% 361_Kagome_FFGFT_Vergleich_De_ch.tex
% Johann Pascher, 12. Sept. 2026
% Strukturelle Parallelen zwischen kinetischer Kagome-Frustration
% (Pei et al., PRL 137, 106702, 2026) und FFGFT

\providecommand{\BM}{\textbf{[B]}}
\providecommand{\KM}{\textbf{[K]}}
\providecommand{\SM}{\textbf{[S]}}
\providecommand{\SetzM}{\textbf{[SETZUNG]}}
\providecommand{\QM}{\textbf{[Q]}}

\chapter*{Kagome-Gitter und FFGFT: Vier strukturelle Parallelen}
\addcontentsline{toc}{chapter}{Kagome-Gitter und FFGFT: Vier strukturelle Parallelen}

\begin{abstract}
Pei, Chen, Castelnovo und Moessner (PRL~137, 106702, 31.~August 2026, Open
Access) untersuchen das einfach dotierte Hubbard-Modell mit unendlicher
Abstoßung auf dem Kagome-Gitter. Sie finden resonante Valenzbindungs-Polaronen
(RVB-Polaronen), störungsfreie Selbstlokalisierung mit Bandmassenanstieg um
sechs Größenordnungen sowie einen Übergang zu $\sqrt{3}\times\sqrt{3}$-Ordnung
im unpolarisierten Grenzfall, den sie am Drei-Zustands-Potts-Modell messen.
Dieses Dokument identifiziert vier strukturelle Parallelen zur FFGFT und belegt
jede mit dem Wortlaut der Quelle. Die Parallelen führen auf dieselbe algebraische
Wurzel: den Frobenius-Automorphismus $\phi: x \mapsto x^3$ der Ordnung~3 auf
endlichen Körpern. Ein Abschnitt benennt, was \emph{nicht} übereinstimmt.
Status aller vier Parallelen: \SM\ (strukturelle Übereinstimmung nachgewiesen,
physikalische Identifikation offen).
\end{abstract}

% ============================================================
\section*{Statusmarkierungen}
% ============================================================
\begin{center}
\begin{tabular}{@{}lp{10cm}@{}}
\textbf{[B]} & Algebraisch bewiesen. \\[3pt]
\textbf{[K]} & Numerisch bestätigt. \\[3pt]
\textbf{[S]} & Strukturell plausibel, physikalische Identifikation offen. \\[3pt]
\textbf{[Q]} & Aus der Quelle (Pei et al.) übernommen, wörtlich zitiert. \\[3pt]
\textbf{[SETZUNG]} & Axiomatische Setzung. \\
\end{tabular}
\end{center}

\noindent Alle wörtlichen Zitate stammen aus Pei et al.\ (2026), Open Access
unter CC-BY; Abschnittsangaben beziehen sich auf die dortige Gliederung.

% ============================================================
\section{Die Quelle: Was Pei et al.\ zeigen}
% ============================================================

\subsection{Modell und Fragestellung}

Das untersuchte Modell ist das Hubbard-Modell mit unendlicher Abstoßung auf
dem Kagome-Gitter, ein einziges Loch im System:
\begin{quote}\QM
\enquote{To gain deeper insight into the role of hole kinetics in determining
magnetism in highly frustrated doped Mott insulators, we consider the
single-hole counter-Nagaoka problem on the kagome lattice, using magnetization
as a tuning parameter.} (Abstract)
\end{quote}
Der Begriff \emph{counter-Nagaoka} ist bedeutsam: Nagaokas Theorem sagt
Ferromagnetismus aus kinetischer Energie eines Lochs auf bipartiten Gittern
voraus. Auf dem Kagome-Gitter ist die Kinetik \emph{frustriert}, und es
entstehen antiferromagnetische Korrelationen:
\begin{quote}\QM
\enquote{When said kinetic energy is frustrated, it was later shown that
antiferromagnetic correlations can also appear.} (Einleitung)
\end{quote}

\subsection{Die Physik des einzelnen Dreiecks}

Der Mechanismus wird an einem einzelnen Dreieck sichtbar. Pei et al.\
beschreiben ihn so:
\begin{quote}\QM
\enquote{Ferromagnetic spins have lowest energy $E_0=-t$ due to destructive
interference of the hole motion around the triangle; whereas antiferromagnetic
spins can form a singlet and effectively thread a $\pi$ flux through the
triangle, relieving the frustration and allowing the hole to delocalize and
lower the energy to $E_0=-2t$ (bottom of an unfrustrated 1D tight-binding
band).} (Abschnitt \emph{Model})
\end{quote}
Zwei Dinge sind hier festzuhalten. Erstens: Die ferromagnetische
Konfiguration ist \emph{destruktiv interferiert} -- das Loch kann sich nicht
frei bewegen. Zweitens: Das Singulett \emph{fädelt einen $\pi$-Fluss} durch
das Dreieck. Dieser $\pi$-Fluss ist keine Metapher, sondern eine Berry-Phase:
das Loch akkumuliert bei einem vollständigen Umlauf die Phase $\pi$, die
Interferenz wechselt von destruktiv zu konstruktiv.

\subsection{Flachband und kompakt lokalisierte Zustände}

Im vollständig polarisierten Fall reduziert sich das Problem auf spinlose
Fermionen. Das unterste Band ist dann exakt flach:
\begin{quote}\QM
\enquote{The fully polarized case reduces to a well-known spinless fermionic
tight-binding model, where the bottom band is exactly flat (energy $-2$) and
the states are compact localized around individual hexagons.}
(Abschnitt \emph{RVB polarons and delocalization})
\end{quote}
Die kinetische Frustration entsteht genau dann, wenn dieses Flachband das
energetisch niedrigste ist:
\begin{quote}\QM
\enquote{Indeed, kinetic frustration arises when the kagome flat band happens
to be lowest in energy for a hole-doped fully polarized electron system close
to half-filling. As such, holes in a fully polarized system occupy compact
localized states at low energy.} (Einleitung)
\end{quote}

\subsection{RVB-Polaronen: Bildung, Delokalisierung, Selbstfalle}

Beim Hinzufügen weniger umgekehrter Spins ($n=1,\ldots,4$) bilden sich
Polaronen, die \emph{delokalisieren} -- aber mit stark erhöhter Bandmasse:
\begin{quote}\QM
\enquote{Upon adding few reversed spins, polarons form and delocalize, albeit
with a band mass 100 times larger than that of an unfrustrated tight-binding
particle on the same lattice.} (Einleitung)
\end{quote}
Die Struktur dieser Polaronen wird durch Valenzbindungssinguletts auf
baumartigen Husimi-Kakteen beschrieben:
\begin{quote}\QM
\enquote{The short distance physics of the kagome lattice is close to
tree-like, and we find that the valence bond language of the same Hamiltonian
on the Husimi cactus proves highly useful to model and understand the polarons
in this Letter.} (Abschnitt \emph{RVB polarons and delocalization})
\end{quote}
Bei $n=5$ ändert sich das Verhalten schlagartig:
\begin{quote}\QM
\enquote{The polaron bandwidth suddenly drops by at least six orders of
magnitude ($<10^{-8}$ in a $4\times4$ system of 48 sites, with 16
near-degenerate GSs); to the best of our numerics it is indistinguishable
from a flat band of localized states.} (Abschnitt \emph{Strong self-trapping})
\end{quote}
Der Mechanismus dieser Selbstfalle ist eine Resonanz sechs
symmetrieverwandter Cluster um eine hexagonale Plakette:
\begin{quote}\QM
\enquote{[\ldots] there are six symmetry-related such clusters around any
hexagonal plaquette of the lattice, and the corresponding six Husimi cactus
states are not orthogonal. [\ldots] the GS wave function takes the form of a
symmetric superposition of all six states, spreading across all 12 sites
around the hexagonal plaquette.} (Abschnitt \emph{Strong self-trapping})
\end{quote}
Und die Deutung:
\begin{quote}\QM
\enquote{By resonating around a hexagonal plaquette, it seems that the 1H5S
RVB polaron acquires a possibly infinite mass, undergoing localization
reminiscent to the compact-localized hexagonal states of the spinless
single-particle problem.} (Abschnitt \emph{Strong self-trapping})
\end{quote}
Bemerkenswert ist, dass diese Lokalisierung in einem \emph{unordnungsfreien}
System eintritt: Die Autoren sprechen von \enquote{relocalize in a
disorder-free system} (Schluss).

\subsection{Der unpolarisierte Grenzfall und das Potts-Modell}

Mit wachsendem $n$ verliert die Valenzbindungssprache ihre Nützlichkeit, und
es entstehen semiklassische Korrelationen:
\begin{quote}\QM
\enquote{We see that this \enquote{frustration on frustration} leads to the
appearance of magnetic correlations consistent with semiclassical
$\sqrt{3}\times\sqrt{3}$ order as we approach the unpolarized limit.}
(Abschnitt \emph{Unpolarized state})
\end{quote}
Als Referenz dient das Drei-Zustands-Potts-Modell:
\begin{quote}\QM
\enquote{The data presented in Fig.~5 include our Letter of the three-state
Potts model, with Hamiltonian $H_\text{Potts}=\sum_{\langle ij\rangle}
\delta_{\sigma_i\sigma_j}$, $\sigma_i=0,1,2$. The ground state ensemble is
extensively degenerate and consists of all states without identical nearest
neighbors.} (Anhang \emph{Potts model})
\end{quote}
Dieses Modell ist nicht willkürlich gewählt -- es markiert einen Phasenübergang:
\begin{quote}\QM
\enquote{This model is useful in that it sits at the Kosterlitz-Thouless (KT)
transition separating a $\sqrt{3}\times\sqrt{3}$ ordered and an algebraically
correlated phase, and it can thus provide a natural benchmark for an analysis
of semiclassical correlations.} (Abschnitt \emph{Unpolarized state})
\end{quote}

% ============================================================
\section{Parallele 1: Frobenius-Orbits -- dieselbe algebraische Struktur}
% ============================================================

\subsection{FFGFT-Seite}

\BM\ (Dok.~339): Die multiplikative Gruppe $GF(27)^*\cong\mathbb{Z}_{26}$
zerfällt unter dem Frobenius-Automorphismus $\phi: x\mapsto x^3$ in
\begin{itemize}
  \item zwei Fixpunkte $\{+1,-1\}$ -- der massive $\mathbb{Z}_2$-Sektor
        (Teilchen/Antiteilchen),
  \item vier Dreier-Orbits in $\mathbb{Z}_{13}$, je $\{a,3a,9a\pmod{13}\}$:
        $\{1,3,9\}$, $\{2,5,6\}$, $\{4,10,12\}$, $\{7,8,11\}$,
  \item kombiniert mit den zwei $\mathbb{Z}_2$-Paritäten: acht Gluonen der
        adjungierten $SU(3)$-Darstellung.
\end{itemize}
Die Dreier-Orbits entsprechen den drei Farbzuständen $\{R,G,B\}$; die
Parität unterscheidet Farbe von Antifarbe (Dok.~339, §3.2).

\subsection{Kagome-Seite}

Das Kagome-Gitter hat drei Untergitterplätze pro Einheitszelle, die
Einheitszelle selbst ist ein Dreieck. Pei et al.\ verwenden explizit
Systemgrößen, die Vielfache von~3 sind:
\begin{quote}\QM
\enquote{$\sqrt{3}\times\sqrt{3}$ order requires $L_x$ and $L_y$ to be
integer multiples of 3 to avoid incommensurability issues.}
(Anhang \emph{Methods})
\end{quote}
Die Ordnung selbst zeigt sich an den K-Punkten der Brillouin-Zone:
\begin{quote}\QM
\enquote{In all cases, we find peaks at the $K$ points, suggestive of
$\sqrt{3}\times\sqrt{3}$ order.} (Abschnitt \emph{Unpolarized state})
\end{quote}
Die K-Punkte tragen dreifache Rotationssymmetrie ($C_3$); der
Strukturfaktor $S(\mathbf{k})$ ist invariant unter der Drehung um $2\pi/3$.

\subsection{Strukturelle Übereinstimmung}

Beide Seiten tragen dieselbe Algebra: eine zyklische Gruppe, auf der ein
Automorphismus der Ordnung~3 wirkt und Dreier-Orbits erzeugt. In FFGFT ist es
$\phi$ auf $GF(27)^*$; auf dem Kagome-Gitter ist es $C_3$ auf den
Untergitterplätzen und -- im reziproken Raum -- auf den K-Punkten. \SM

Was hier \emph{nicht} behauptet wird: dass die vier $\mathbb{Z}_{13}$-Orbits
physikalisch mit vier Kagome-Konfigurationen identisch sind. Das Kagome-Gitter
hat einen einzigen $C_3$-Orbit pro Einheitszelle; FFGFT hat vier Orbits in
$\mathbb{Z}_{13}$. Die Übereinstimmung liegt in der \emph{Ordnung} des
wirkenden Automorphismus, nicht in der Anzahl der Orbits.

% ============================================================
\section{Parallele 2: Drei-Zustands-Potts-Modell = $GF(3)$ = Fixkörper}
% ============================================================

\subsection{FFGFT-Seite}

\BM\ (Dok.~339, §3.3): Der Fixkörper von $\phi$ auf $GF(27)$ ist
$GF(3)=\{0,+1,-1\}$ (kleiner Fermat: $x^3=x$ für alle $x\in GF(3)$). In
FFGFT entspricht $GF(3)$ dem $U(1)$-Sektor: die Projektion $GF(27)\to GF(3)$
beschreibt den elektrisch neutralen Sektor (Photon). Die drei Elemente
$\{0,+1,-1\}$ sind zugleich die drei $\mathbb{Z}_3$-Sektoren aus Dok.~336
(Fixpunkte $\{0,+1,-1\}$ = FFGFT-Sektoren \BM).

\BM\ (Dok.~321, §1.2): Die Triality-Bedingung $T_R=0\pmod{3}$ ist
äquivalent zu Confinement:
\begin{quote}
\enquote{Confinement is equivalent to requiring all observable states to have
triality $T_R=0$.} (Dok.~321)
\end{quote}
Ein einzelnes Quark ($T_R=1$) ist nicht observabel; Meson ($T_R=1+2=0$) und
Baryon ($T_R=1+1+1=0$) sind es.

\subsection{Kagome-Seite}

Das Potts-Benchmark-Modell hat exakt drei Zustände $\sigma\in\{0,1,2\}$ mit
$\mathbb{Z}_3$-Symmetrie (siehe Zitat oben). Die Grundzustandsbedingung
\enquote{no identical nearest neighbors} ist eine Frustrationsbedingung: auf
einem Dreieck müssen alle drei Zustände vorkommen -- genau eine
$\mathbb{Z}_3$-neutrale Konfiguration pro Dreieck. Der KT-Übergang trennt die
Phase, in der diese $\mathbb{Z}_3$-Ordnung langreichweitig ist
($\sqrt{3}\times\sqrt{3}$), von der algebraisch korrelierten Phase.

\subsection{Strukturelle Übereinstimmung}

Die Identifikation $\sigma\in\{0,1,2\}\leftrightarrow GF(3)=\{0,+1,-1\}$ ist
algebraisch exakt: beide sind $\mathbb{Z}_3$ als additive Gruppe. Die
Frustrationsbedingung des Potts-Modells (alle drei Zustände auf jedem
Dreieck) und die Confinement-Bedingung der FFGFT (Summe der Trialitäten
$\equiv0\pmod{3}$) sind dieselbe Aussage: nur $\mathbb{Z}_3$-neutrale
Konfigurationen sind zulässig. \SM

Der KT-Übergang des Potts-Modells hat damit ein FFGFT-Gegenstück: den
Übergang zwischen confined ($T_R=0$ langreichweitig durchgesetzt) und
deconfined (Triality nur algebraisch korreliert). Ob diese Identifikation
über die Algebra hinaus physikalisch trägt, ist offen (§7).

% ============================================================
\section{Parallele 3: Masse = eingefrorene kinetische Energie}
% ============================================================

\subsection{FFGFT-Seite}

\BM\ (Dok.~286): Die Dualität $\tilde{T}\cdot m=1$ erzwingt zwei
Energietypen:
\begin{center}
\begin{tabular}{lll}
\toprule
 & Statischer Sektor & Kinetischer Sektor \\
\midrule
Form        & $E_0=mc^2$ (Ruhemasse)         & $E=hf$ (frei) \\
Sitz        & $\mathbb{Z}_3$-Kern (Wicklung) & Fenster \\
Zeit        & $\tilde{T}=1/m$                & $\tilde{T}=1/\omega$ \\
Status      & hergeleitet (Dok.~285)         & offen (Dok.~286) \\
\bottomrule
\end{tabular}
\end{center}
Dok.~286 formuliert: \enquote{The rest mass is static: a topologically stable
winding number on $T^4$ [\ldots] \emph{Mass is frozen kinetic energy.}}
\BM\ (Dok.~339): Die Trennung massiv/masselos ist durch den Frobenius
algebraisch erzwungen -- Fixpunkte $\{+1,-1\}$ sind massiv, die
Dreier-Orbits sind masselos (Gluonen).

\subsection{Kagome-Seite}

Das Kagome-System zeigt \emph{dieselbe Zweiteilung}, aber nicht als
Sektor-Trennung, sondern als \emph{nichtmonotone Bandmasse} beim Durchlaufen
der Polarisation:
\begin{itemize}
  \item $n=0$: Flachband, kompakt lokalisiert -- \enquote{effektiv unendliche
        Masse} (kinetische Energie eingefroren durch destruktive Interferenz).
  \item $n=1,\ldots,4$: Delokalisierung mit Bandmasse $\times100$
        (\enquote{polarons form and delocalize}).
  \item $n=5$: Relokalisierung, Bandmasse $\times10^6$ (\enquote{acquires a
        possibly infinite mass}).
  \item $n\to$ unpolarisiert: Delokalisierung (\enquote{delocalized, dispersive
        holes are recovered}).
\end{itemize}
Pei et al.\ fassen zusammen:
\begin{quote}\QM
\enquote{Our Letter uncovers an intricate interplay between single-particle
flat bands and cooperative \enquote{polaronic} self-trapping physics.}
(Einleitung)
\end{quote}

\subsection{Strukturelle Übereinstimmung}

In beiden Systemen ist \emph{Masse} das Ergebnis geometrischer Einfrierung
kinetischer Freiheit: in FFGFT durch topologische Wicklung auf $T^4$
(Dok.~314: \enquote{winding numbers are topological}), auf dem Kagome-Gitter
durch destruktive Interferenz um Plaketten. Beides ist Lokalisierung
\emph{ohne Unordnung} -- ein rein geometrischer Effekt. \SM

Der Unterschied: FFGFT hat zwei \emph{Sektoren} (massiv/masselos), das
Kagome-System einen \emph{Parameter} (Polarisation), entlang dessen die
Bandmasse zwischen beiden Extremen wechselt. Das Kagome-System realisiert
damit den Übergang, den FFGFT als Sektorgrenze behandelt.

% ============================================================
\section{Parallele 4: $\pi$-Fluss = Berry-Phase = Wicklungszahl}
% ============================================================

\subsection{Kagome-Seite}

Das Zitat aus §1.2 ist hier zentral: das Singulett \enquote{effectively
thread[s] a $\pi$ flux through the triangle}. Der $\pi$-Fluss ist die Phase,
die das Loch bei einem Umlauf um das Dreieck aufnimmt. Ohne diese Phase
(ferromagnetisch) interferiert die Hüpfamplitude destruktiv; mit ihr
konstruktiv. Der Energiegewinn ist ein Faktor~2 ($E_0=-t\to-2t$).

Dieselbe Logik wiederholt sich auf größerer Skala bei der Selbstfalle: die
sechs Husimi-Zustände um die Plakette interferieren \emph{symmetrisch} --
das ist konstruktive Interferenz für die gebundene Komponente, die den
Zustand um alle zwölf Plätze verteilt, aber destruktive für jede
Bewegung \emph{aus} der Plakette heraus.

\subsection{FFGFT-Seite}

\BM\ (Dok.~314, Ch.~D2): \enquote{in the rolled-up state there are no
fractal corrections -- winding numbers are topological, the defect $d=100\xi$
accumulates only when unrolled.} Wicklungszahlen auf $T^4$ sind quantisierte
Berry-Phasen um geschlossene Schleifen: die Phase $2\pi n$, $n\in\mathbb{Z}$,
bei einem Umlauf um einen Zyklus des Torus.

\BM\ (Dok.~321, §2.1): Der $\mathbb{Z}_3$-Generator $\tau$ wirkt als unitärer
Operator auf $L^2(T^4)$ mit $\tau^3=\mathrm{id}$; die Eigenwerte sind
$\omega^k$, $\omega=e^{2\pi i/3}$. Die drei Sektoren $P_k$ sind durch die
akkumulierte Phase bei einer $\mathbb{Z}_3$-Drehung unterschieden.

\BM\ (Dok.~230): Verschränkung ist \enquote{Holonomie um Zyklen} -- die
akkumulierte Phase bei einem geschlossenen Umlauf ist das, was
Verschränkung geometrisch ausmacht.

\subsection{Strukturelle Übereinstimmung}

Der $\pi$-Fluss im Kagome-Dreieck (Berry-Phase $\pi$, halbe Wicklung) und
die $\mathbb{Z}_3$-Phasen $\omega^k$ in FFGFT (Berry-Phase $2\pi k/3$) sind
dieselbe Operation: akkumulierte Phase entlang eines geschlossenen Pfads, die
den zugänglichen Hilbertraum verändert. Auf dem Kagome-Gitter hebt sie die
Frustration auf; in FFGFT selektiert sie den Sektor. \SM

% ============================================================
\section{Was nicht übereinstimmt}
% ============================================================

Ehrlichkeit verlangt, die Grenzen der Parallelen zu benennen.

\paragraph{Dimension.} Das Kagome-Gitter ist zweidimensional; $T^4/Z_3$ ist
vierdimensional. Die $C_3$-Symmetrie des Kagome-Gitters lebt in der Ebene;
die $\mathbb{Z}_3$-Wirkung auf $T^4$ hat neun Fixpunkte (Dok.~314) und
wirkt auf einem kompakten Raum ohne Rand. Es gibt keine Abbildung, die das
Kagome-Gitter in $T^4/Z_3$ einbettet.

\paragraph{Klassisch versus quantenmechanisch.} Das Potts-Modell ist
klassisch; Pei et al.\ verwenden es als \emph{semiklassische} Referenz und
skalieren Korrelatoren mit $(S^z)^2=1/4$ statt $S(S+1)=3/4$. FFGFT arbeitet
auf $L^2(T^4)\otimes\mathbb{C}^3$ (Dok.~230/282) durchgehend
quantenmechanisch. Die Übereinstimmung in Parallele~2 ist algebraisch
($\mathbb{Z}_3$), nicht dynamisch.

\paragraph{Anzahl der Orbits.} FFGFT hat vier Dreier-Orbits in
$\mathbb{Z}_{13}$ (acht Gluonen); das Kagome-Gitter hat einen $C_3$-Orbit
pro Einheitszelle. Parallele~1 betrifft die \emph{Ordnung} des Automorphismus,
nicht die Orbitzahl.

\paragraph{Struktur versus Dynamik.} Alle vier Parallelen sind strukturell.
Keine sagt voraus, dass ein FFGFT-Rechenweg das Kagome-Spektrum
reproduziert oder umgekehrt. Sie sagen: dieselbe Algebra ($\phi$ der
Ordnung~3) erzeugt in beiden Systemen dieselbe Art von Einschränkung
(nur $\mathbb{Z}_3$-neutrale Konfigurationen zulässig; Lokalisierung ohne
Unordnung; Berry-Phase als Selektor).

\paragraph{Eine unaufgeklärte Zahl.} Der Selbstfallen-Cluster umfasst zwölf
Plätze um eine hexagonale Plakette. In Dok.~314 hat die erste Anregung des
$D_4$-Gitters Entartung $24=12+12$ (\enquote{FCC half and winding half
separately addressable}). Ob die~12 des Kagome-Clusters mit einer der
12-Hälften von $D_4$ zusammenhängt, ist \emph{nicht} untersucht und wird hier
nicht als Parallele geführt. \SM

% ============================================================
\section{Zusammenfassung und offene Fragen}
% ============================================================

\begin{center}
\begin{tabular}{p{4.2cm}p{4.2cm}p{3.6cm}l}
\toprule
Kagome-System (Pei et al.) & FFGFT & Gemeinsame Algebra & Status \\
\midrule
K-Punkt-Ordnung, $C_3$ auf Untergitter &
  Frobenius-Orbits auf $GF(27)^*$ (Dok.~339) &
  Automorphismus Ordnung~3 & \SM \\
Drei-Zustands-Potts, KT-Übergang &
  $GF(3)$ Fixkörper, Triality $T_R=0$ (Dok.~339, 321) &
  $\mathbb{Z}_3$-Neutralität & \SM \\
Flachband / Selbstfalle, Bandmasse $\times10^6$ &
  Statisch/kinetisch, eingefrorene Wicklung (Dok.~286, 314) &
  Lokalisierung ohne Unordnung & \SM \\
$\pi$-Fluss durch Dreieck &
  Wicklungszahl, $\mathbb{Z}_3$-Phase $\omega^k$ (Dok.~314, 321) &
  Berry-Phase um Schleife & \SM \\
\bottomrule
\end{tabular}
\end{center}

Alle vier Parallelen sind strukturell nachgewiesen \SM. Die physikalische
Frage -- ob das Kagome-System eine Niederenergie-Realisierung der
FFGFT-Galois-Struktur ist -- bleibt offen und wird hier nicht behauptet.

\paragraph{Offene Fragen \SM.}
\begin{enumerate}
  \item Ist der KT-Übergang des Potts-Modells (Trennung
        $\mathbb{Z}_3$-Ordnung / algebraisch korreliert) mit dem
        FFGFT-Confinement-Übergang ($T_R=0$ durchgesetzt / nur korreliert)
        bei endlicher Temperatur identifizierbar?
  \item Pei et al.\ kündigen für zwei Löcher \enquote{polaron-mediated
        interactions [\ldots] and possible binding, pairing} an. Zwei Löcher
        mit je $T_R\neq0$ könnten in FFGFT-Sprache einen $T_R=0$-Zustand bilden.
        Ist das Loch-Paar-Binding auf dem Kagome-Gitter ein Meson-Analogon?
  \item Die zwölf Plätze des Selbstfallen-Clusters gegen $24=12+12$ aus
        Dok.~314 -- Zufall oder Struktur?
\end{enumerate}

\paragraph{Quelldokument.} Pei,~Y.; Chen,~S.~A.; Castelnovo,~C.;
Moessner,~R.: \textit{Kinetic Kagome Magnetism: From Self-Trapping RVB
Polarons to Semiclassical Correlations.}
Phys.~Rev.~Lett.~\textbf{137}, 106702, veröffentlicht 31.~August 2026.
DOI: 10.1103/zbxv-vtq8. Open Access (Editors' Suggestion).
